\section{系统架构与数据流}

\subsection{整体架构}

Insight 金融研报分析引擎采用高内聚、低耦合的模块化设计，完美适配 Mac Studio (M3 Ultra)
的极致本地算力。系统由四个核心层次组成：数据接入层、检索生成层、评估闭环层和交互层。

\subsection{系统架构图说明}

系统架构图包含以下核心组件：

1. **数据接入层 (Ingest)**：负责 PDF 解析、分块和向量化
2. **检索生成层 (Retrieval \& Generation)**：负责信息检索和回答生成
3. **评估闭环层 (Evaluation)**：负责系统性能评估
4. **交互层 (UI)**：负责用户界面和交互

数据流向：PDF 研报 → 解析 → 分块 → 向量化 → 存储 → 检索 → 精排 → 生成回答 → 评估 → 展示

\subsection{数据接入层 (Ingest)}

数据接入层负责金融 PDF 研报的解析、分块和向量化，为后续的检索和生成提供基础数据。

\subsubsection{PDF 解析}
使用 PyMuPDFReader 对 PDF 研报进行解析，提取文本内容并保留页码信息。系统确保每个文档的元数据中包含 `source`（文件名）和 `page\_label`（物理页码），这是实现引用溯源的基础。

\subsubsection{分块策略}
系统支持两种分块策略：

1. **基线分块 (Phase 1)**：使用固定大小的分块策略，块大小为 256 tokens，重叠度为 10%。
2. **语义分块 (Phase 2)**：基于 BGE-M3 的语义分块，通过计算句间余弦相似度，在语义发生突变时进行切分，最大程度保留财务分析的上下文完整性。

\subsubsection{向量化与存储}
使用 BAAI/bge-m3 模型对分块后的文本进行向量化，生成密集向量，并存储到本地 ChromaDB 向量库中。对于 Phase 2，系统还会构建 RAPTOR 树状摘要，利用后台小模型对底层语义块进行递归聚类与摘要，构建宏观知识树。

\subsection{检索生成层 (Retrieval \& Generation)}

检索生成层负责根据用户查询从向量库中检索相关信息，并通过大语言模型生成回答。

\subsubsection{检索策略}
系统支持两种检索策略：

1. **单路稠密向量检索 (Phase 1)**：直接使用 VectorStoreIndex 进行相似度检索。
2. **多路混合检索 (Phase 2)**：集成 BM25 检索和向量检索，通过 RRF（Reciprocal Rank Fusion）融合结果，并使用 bge-reranker-base 进行精排。

\subsubsection{对话引擎}
系统使用 CondensePlusContextChatEngine 作为核心对话引擎，集成 ChatMemoryBuffer 实现短期记忆管理，确保模型能在多轮对话中理解上下文指代。

\subsubsection{生成策略}
系统使用本地部署的 Qwen2.5-7B 或 Qwen2.5-72B 模型生成回答。在生成过程中，系统强制要求输出引用信息，确保回答的可验证性。

\subsection{评估闭环层 (Evaluation)}

评估闭环层负责对系统的性能进行评估，确保系统输出的质量。

系统使用 RAGAS 评估框架，将本地运行的 Qwen2.5-72B 注册为 RAGAS 的 judge_model，对系统回答进行打分，评估指标包括：

1. **Context Precision**：衡量系统召回的文档块中，真正包含答案相关信息的块是否排在最前面。
2. **Faithfulness**：衡量大模型生成的回答是否 100% 由召回的文档块推导得出，有无幻觉。
3. **Answer Relevancy**：衡量回答是否直接命中了用户的原始问题。

\subsection{交互层 (UI)}

交互层负责构建用户界面，提供友好的用户体验。

系统使用 Gradio 构建前端界面，实现以下功能：

1. 多文档管理：支持上传和管理多个 PDF 研报。
2. 对话式问答：提供自然语言对话界面，支持多轮对话。
3. 引用跳转：实现点击引用跳转到 PDF 对应页码的功能。
4. 产物生成：支持生成执行摘要、FAQ 等结构化产物。

\subsection{数据流}

系统的端到端数据流如下：

1. **数据接入**：PDF 研报 → PyMuPDF 解析 → 分块 → BGE-M3 向量化 → ChromaDB 存储
2. **检索生成**：用户查询 → 检索策略分发 → 混合检索 → 精排 → 大模型生成回答
3. **评估**：系统回答 → RAGAS 评估 → 量化指标计算
4. **交互**：用户输入 → 前端处理 → 后端处理 → 前端展示
